\documentclass[11pt,a4paper]{article}
% lesson-preamble.tex -- Shared preamble for all Phase 00 lessons
% Usage: % lesson-preamble.tex -- Shared preamble for all Phase 00 lessons
% Usage: % lesson-preamble.tex -- Shared preamble for all Phase 00 lessons
% Usage: \input{../lesson-preamble} (from subdomain directories)

% ---- Encoding and fonts ----
\usepackage[utf8]{inputenc}
\usepackage[T1]{fontenc}

% ---- Mathematics ----
\usepackage{amsmath,amssymb,amsthm}
\usepackage{mathtools}

% ---- Page geometry ----
\usepackage[margin=1in,headheight=14pt]{geometry}

% ---- Hyperlinks ----
\usepackage[colorlinks=true,linkcolor=red,citecolor=blue,urlcolor=blue]{hyperref}

% ---- Lists ----
\usepackage{enumitem}

% ---- Colors and boxes ----
\usepackage{xcolor}
\usepackage[theorems,skins,breakable]{tcolorbox}

% ---- Header/Footer ----
\usepackage{fancyhdr}
\renewcommand{\headrulewidth}{0.4pt}

% ---- Tables ----
\usepackage{booktabs}

% ---- Bibliography ----
\usepackage{natbib}
\bibliographystyle{plainnat}

% --------------------------------------------------------------------------
% Theorem environments
% --------------------------------------------------------------------------
\theoremstyle{plain}
\newtheorem{theorem}{Theorem}[section]
\newtheorem{lemma}[theorem]{Lemma}
\newtheorem{proposition}[theorem]{Proposition}
\newtheorem{corollary}[theorem]{Corollary}

\theoremstyle{definition}
\newtheorem{definition}[theorem]{Definition}
\newtheorem{example}[theorem]{Example}
\newtheorem{exercise}[theorem]{Exercise}
\newtheorem{assumption}[theorem]{Assumption}

\theoremstyle{remark}
\newtheorem{remark}[theorem]{Remark}

% --------------------------------------------------------------------------
% Colored boxes
% --------------------------------------------------------------------------
\newtcolorbox{keyresult}[1][]{%
  colback=green!5!white,
  colframe=green!50!black,
  fonttitle=\bfseries,
  title={Key Result},
  breakable,
  #1
}

\newtcolorbox{rlconnection}[1][]{%
  colback=blue!5!white,
  colframe=blue!50!black,
  fonttitle=\bfseries,
  title={RL/ML Connection},
  breakable,
  #1
}

\newtcolorbox{warning}[1][]{%
  colback=red!5!white,
  colframe=red!50!black,
  fonttitle=\bfseries,
  title={Warning},
  breakable,
  #1
}

% --------------------------------------------------------------------------
% Custom commands -- number sets
% --------------------------------------------------------------------------
\newcommand{\R}{\mathbb{R}}
\newcommand{\N}{\mathbb{N}}
\newcommand{\Q}{\mathbb{Q}}
\newcommand{\Z}{\mathbb{Z}}
\newcommand{\C}{\mathbb{C}}
\newcommand{\F}{\mathbb{F}}

% --------------------------------------------------------------------------
% Custom commands -- probability and expectation
% --------------------------------------------------------------------------
\newcommand{\E}{\mathbb{E}}
\newcommand{\Prob}{\mathbb{P}}
\newcommand{\Var}{\operatorname{Var}}
\newcommand{\Cov}{\operatorname{Cov}}
\newcommand{\indicator}{\mathbf{1}}

% --------------------------------------------------------------------------
% Custom commands -- convergence arrows
% --------------------------------------------------------------------------
\newcommand{\cas}{\xrightarrow{\text{a.s.}}}
\newcommand{\cprob}{\xrightarrow{\Prob}}
\newcommand{\clp}[1]{\xrightarrow{L^{#1}}}
\newcommand{\cdist}{\xrightarrow{d}}

% --------------------------------------------------------------------------
% Custom commands -- common operators
% --------------------------------------------------------------------------
\DeclareMathOperator{\dom}{dom}
\DeclareMathOperator{\epi}{epi}
\DeclareMathOperator{\conv}{conv}
\DeclareMathOperator{\relint}{relint}
\DeclareMathOperator{\aff}{aff}
\DeclareMathOperator{\cl}{cl}
\DeclareMathOperator{\interior}{int}
\DeclareMathOperator{\tr}{tr}
\DeclareMathOperator{\diag}{diag}
\DeclareMathOperator{\rank}{rank}
\DeclareMathOperator{\sgn}{sgn}
\DeclareMathOperator{\supp}{supp}
\DeclareMathOperator{\argmin}{arg\,min}
\DeclareMathOperator{\argmax}{arg\,max}
\DeclareMathOperator{\prox}{prox}
\DeclareMathOperator{\dist}{dist}
\DeclareMathOperator{\diam}{diam}
\DeclareMathOperator{\Span}{span}

% --------------------------------------------------------------------------
% Custom commands -- linear algebra operators
% --------------------------------------------------------------------------
\DeclareMathOperator{\nullity}{nullity}
\DeclareMathOperator{\range}{range}
\DeclareMathOperator{\nullspace}{null}
\DeclareMathOperator{\proj}{proj}

% --------------------------------------------------------------------------
% Custom commands -- measure theory
% --------------------------------------------------------------------------
\newcommand{\powerset}{\mathcal{P}}
\newcommand{\Borel}{\mathcal{B}}
 (from subdomain directories)

% ---- Encoding and fonts ----
\usepackage[utf8]{inputenc}
\usepackage[T1]{fontenc}

% ---- Mathematics ----
\usepackage{amsmath,amssymb,amsthm}
\usepackage{mathtools}

% ---- Page geometry ----
\usepackage[margin=1in,headheight=14pt]{geometry}

% ---- Hyperlinks ----
\usepackage[colorlinks=true,linkcolor=red,citecolor=blue,urlcolor=blue]{hyperref}

% ---- Lists ----
\usepackage{enumitem}

% ---- Colors and boxes ----
\usepackage{xcolor}
\usepackage[theorems,skins,breakable]{tcolorbox}

% ---- Header/Footer ----
\usepackage{fancyhdr}
\renewcommand{\headrulewidth}{0.4pt}

% ---- Tables ----
\usepackage{booktabs}

% ---- Bibliography ----
\usepackage{natbib}
\bibliographystyle{plainnat}

% --------------------------------------------------------------------------
% Theorem environments
% --------------------------------------------------------------------------
\theoremstyle{plain}
\newtheorem{theorem}{Theorem}[section]
\newtheorem{lemma}[theorem]{Lemma}
\newtheorem{proposition}[theorem]{Proposition}
\newtheorem{corollary}[theorem]{Corollary}

\theoremstyle{definition}
\newtheorem{definition}[theorem]{Definition}
\newtheorem{example}[theorem]{Example}
\newtheorem{exercise}[theorem]{Exercise}
\newtheorem{assumption}[theorem]{Assumption}

\theoremstyle{remark}
\newtheorem{remark}[theorem]{Remark}

% --------------------------------------------------------------------------
% Colored boxes
% --------------------------------------------------------------------------
\newtcolorbox{keyresult}[1][]{%
  colback=green!5!white,
  colframe=green!50!black,
  fonttitle=\bfseries,
  title={Key Result},
  breakable,
  #1
}

\newtcolorbox{rlconnection}[1][]{%
  colback=blue!5!white,
  colframe=blue!50!black,
  fonttitle=\bfseries,
  title={RL/ML Connection},
  breakable,
  #1
}

\newtcolorbox{warning}[1][]{%
  colback=red!5!white,
  colframe=red!50!black,
  fonttitle=\bfseries,
  title={Warning},
  breakable,
  #1
}

% --------------------------------------------------------------------------
% Custom commands -- number sets
% --------------------------------------------------------------------------
\newcommand{\R}{\mathbb{R}}
\newcommand{\N}{\mathbb{N}}
\newcommand{\Q}{\mathbb{Q}}
\newcommand{\Z}{\mathbb{Z}}
\newcommand{\C}{\mathbb{C}}
\newcommand{\F}{\mathbb{F}}

% --------------------------------------------------------------------------
% Custom commands -- probability and expectation
% --------------------------------------------------------------------------
\newcommand{\E}{\mathbb{E}}
\newcommand{\Prob}{\mathbb{P}}
\newcommand{\Var}{\operatorname{Var}}
\newcommand{\Cov}{\operatorname{Cov}}
\newcommand{\indicator}{\mathbf{1}}

% --------------------------------------------------------------------------
% Custom commands -- convergence arrows
% --------------------------------------------------------------------------
\newcommand{\cas}{\xrightarrow{\text{a.s.}}}
\newcommand{\cprob}{\xrightarrow{\Prob}}
\newcommand{\clp}[1]{\xrightarrow{L^{#1}}}
\newcommand{\cdist}{\xrightarrow{d}}

% --------------------------------------------------------------------------
% Custom commands -- common operators
% --------------------------------------------------------------------------
\DeclareMathOperator{\dom}{dom}
\DeclareMathOperator{\epi}{epi}
\DeclareMathOperator{\conv}{conv}
\DeclareMathOperator{\relint}{relint}
\DeclareMathOperator{\aff}{aff}
\DeclareMathOperator{\cl}{cl}
\DeclareMathOperator{\interior}{int}
\DeclareMathOperator{\tr}{tr}
\DeclareMathOperator{\diag}{diag}
\DeclareMathOperator{\rank}{rank}
\DeclareMathOperator{\sgn}{sgn}
\DeclareMathOperator{\supp}{supp}
\DeclareMathOperator{\argmin}{arg\,min}
\DeclareMathOperator{\argmax}{arg\,max}
\DeclareMathOperator{\prox}{prox}
\DeclareMathOperator{\dist}{dist}
\DeclareMathOperator{\diam}{diam}
\DeclareMathOperator{\Span}{span}

% --------------------------------------------------------------------------
% Custom commands -- linear algebra operators
% --------------------------------------------------------------------------
\DeclareMathOperator{\nullity}{nullity}
\DeclareMathOperator{\range}{range}
\DeclareMathOperator{\nullspace}{null}
\DeclareMathOperator{\proj}{proj}

% --------------------------------------------------------------------------
% Custom commands -- measure theory
% --------------------------------------------------------------------------
\newcommand{\powerset}{\mathcal{P}}
\newcommand{\Borel}{\mathcal{B}}
 (from subdomain directories)

% ---- Encoding and fonts ----
\usepackage[utf8]{inputenc}
\usepackage[T1]{fontenc}

% ---- Mathematics ----
\usepackage{amsmath,amssymb,amsthm}
\usepackage{mathtools}

% ---- Page geometry ----
\usepackage[margin=1in,headheight=14pt]{geometry}

% ---- Hyperlinks ----
\usepackage[colorlinks=true,linkcolor=red,citecolor=blue,urlcolor=blue]{hyperref}

% ---- Lists ----
\usepackage{enumitem}

% ---- Colors and boxes ----
\usepackage{xcolor}
\usepackage[theorems,skins,breakable]{tcolorbox}

% ---- Header/Footer ----
\usepackage{fancyhdr}
\renewcommand{\headrulewidth}{0.4pt}

% ---- Tables ----
\usepackage{booktabs}

% ---- Bibliography ----
\usepackage{natbib}
\bibliographystyle{plainnat}

% --------------------------------------------------------------------------
% Theorem environments
% --------------------------------------------------------------------------
\theoremstyle{plain}
\newtheorem{theorem}{Theorem}[section]
\newtheorem{lemma}[theorem]{Lemma}
\newtheorem{proposition}[theorem]{Proposition}
\newtheorem{corollary}[theorem]{Corollary}

\theoremstyle{definition}
\newtheorem{definition}[theorem]{Definition}
\newtheorem{example}[theorem]{Example}
\newtheorem{exercise}[theorem]{Exercise}
\newtheorem{assumption}[theorem]{Assumption}

\theoremstyle{remark}
\newtheorem{remark}[theorem]{Remark}

% --------------------------------------------------------------------------
% Colored boxes
% --------------------------------------------------------------------------
\newtcolorbox{keyresult}[1][]{%
  colback=green!5!white,
  colframe=green!50!black,
  fonttitle=\bfseries,
  title={Key Result},
  breakable,
  #1
}

\newtcolorbox{rlconnection}[1][]{%
  colback=blue!5!white,
  colframe=blue!50!black,
  fonttitle=\bfseries,
  title={RL/ML Connection},
  breakable,
  #1
}

\newtcolorbox{warning}[1][]{%
  colback=red!5!white,
  colframe=red!50!black,
  fonttitle=\bfseries,
  title={Warning},
  breakable,
  #1
}

% --------------------------------------------------------------------------
% Custom commands -- number sets
% --------------------------------------------------------------------------
\newcommand{\R}{\mathbb{R}}
\newcommand{\N}{\mathbb{N}}
\newcommand{\Q}{\mathbb{Q}}
\newcommand{\Z}{\mathbb{Z}}
\newcommand{\C}{\mathbb{C}}
\newcommand{\F}{\mathbb{F}}

% --------------------------------------------------------------------------
% Custom commands -- probability and expectation
% --------------------------------------------------------------------------
\newcommand{\E}{\mathbb{E}}
\newcommand{\Prob}{\mathbb{P}}
\newcommand{\Var}{\operatorname{Var}}
\newcommand{\Cov}{\operatorname{Cov}}
\newcommand{\indicator}{\mathbf{1}}

% --------------------------------------------------------------------------
% Custom commands -- convergence arrows
% --------------------------------------------------------------------------
\newcommand{\cas}{\xrightarrow{\text{a.s.}}}
\newcommand{\cprob}{\xrightarrow{\Prob}}
\newcommand{\clp}[1]{\xrightarrow{L^{#1}}}
\newcommand{\cdist}{\xrightarrow{d}}

% --------------------------------------------------------------------------
% Custom commands -- common operators
% --------------------------------------------------------------------------
\DeclareMathOperator{\dom}{dom}
\DeclareMathOperator{\epi}{epi}
\DeclareMathOperator{\conv}{conv}
\DeclareMathOperator{\relint}{relint}
\DeclareMathOperator{\aff}{aff}
\DeclareMathOperator{\cl}{cl}
\DeclareMathOperator{\interior}{int}
\DeclareMathOperator{\tr}{tr}
\DeclareMathOperator{\diag}{diag}
\DeclareMathOperator{\rank}{rank}
\DeclareMathOperator{\sgn}{sgn}
\DeclareMathOperator{\supp}{supp}
\DeclareMathOperator{\argmin}{arg\,min}
\DeclareMathOperator{\argmax}{arg\,max}
\DeclareMathOperator{\prox}{prox}
\DeclareMathOperator{\dist}{dist}
\DeclareMathOperator{\diam}{diam}
\DeclareMathOperator{\Span}{span}

% --------------------------------------------------------------------------
% Custom commands -- linear algebra operators
% --------------------------------------------------------------------------
\DeclareMathOperator{\nullity}{nullity}
\DeclareMathOperator{\range}{range}
\DeclareMathOperator{\nullspace}{null}
\DeclareMathOperator{\proj}{proj}

% --------------------------------------------------------------------------
% Custom commands -- measure theory
% --------------------------------------------------------------------------
\newcommand{\powerset}{\mathcal{P}}
\newcommand{\Borel}{\mathcal{B}}

\pagestyle{fancy}
\fancyhf{}
\fancyhead[L]{\textsc{Lesson 5: Vector Spaces and Linear Maps}}
\fancyhead[R]{\thepage}
\title{Lesson 5: Vector Spaces and Linear Maps}
\author{AI/ML Mathematics}
\date{}
\begin{document}
\maketitle
\tableofcontents
\newpage

%% ============================================================
%% SECTION 1: MOTIVATION AND OVERVIEW
%% ============================================================
\section{Motivation and Overview}
\label{sec:motivation}

Consider a neural network with a single hidden layer.  The forward pass computes
$x \mapsto \sigma(W_2 \,\sigma(W_1 x + b_1) + b_2)$, where $W_1 \in \R^{d \times n}$
and $W_2 \in \R^{k \times d}$ are weight matrices.  Strip away the nonlinearity
$\sigma$ for a moment: the remaining map $x \mapsto W_2 W_1 x$ is a \emph{composition
of linear maps} $\R^n \to \R^d \to \R^k$.  Its range, null space, and rank determine
which inputs the network can distinguish and how many independent features the hidden
layer can encode.  If $\rank(W_2 W_1) < \min(k, n)$, the network suffers a
\emph{bottleneck} -- it collapses distinct inputs to the same output.

More broadly, policy gradient methods in reinforcement learning parameterize a policy
$\pi_\theta$ by a vector $\theta \in \R^p$.  The parameter space is a vector space, and
the Fisher information matrix defines a natural inner product on its tangent space.
Understanding bases, dimension, and change of basis is prerequisite to reasoning about
such geometric structures.  The tools developed in this lesson will let us define linear
maps rigorously, prove the rank--nullity theorem, and analyze how matrix representations
change with a choice of basis -- all of which are essential for the spectral and
optimization theory developed in later lessons.

\begin{rlconnection}[title={Where Vector Spaces and Linear Maps Appear in RL/ML}]
\begin{itemize}[nosep]
  \item \textbf{Linear function approximation (TD learning).}
        Value functions are approximated as $\hat{V}(s) = \phi(s)^\top w$, a linear
        map from feature space to $\R$; the rank of the feature matrix $\Phi$
        determines expressiveness.
  \item \textbf{Policy gradient / natural gradient.}
        The parameter update $\theta \leftarrow \theta + \alpha \, F^{-1} \nabla J$
        requires understanding coordinate changes and similarity transformations.
  \item \textbf{Principal component analysis (PCA).}
        Dimensionality reduction projects data onto a subspace spanned by the top
        eigenvectors -- a direct application of bases, span, and projections.
\end{itemize}
\end{rlconnection}

\paragraph{Prerequisites.}
Familiarity with sets, functions, and the field axioms of $\R$ (Lesson 1);
sequences and limits in $\R$ (Lesson 2); continuity of real-valued functions
(Lesson 3); basic properties of $\R^n$ as a set of tuples.

\paragraph{Learning objectives.}
After completing this lesson, the student will be able to:
\begin{itemize}[nosep]
  \item \textbf{Define} a vector space, subspace, span, linear independence, basis,
        dimension, and linear map, and verify these definitions in concrete examples.
  \item \textbf{Prove} that the intersection of subspaces is a subspace and that any
        two bases of a finite-dimensional space have the same cardinality.
  \item \textbf{Derive} the rank--nullity theorem from the definition of a basis
        and apply it to determine injectivity and surjectivity of linear maps.
  \item \textbf{Compute} the matrix representation of a linear map with respect to
        given bases and carry out a change-of-basis computation.
  \item \textbf{Classify} linear maps between finite-dimensional spaces as injective,
        surjective, or bijective using dimension arguments.
\end{itemize}

%% ============================================================
%% SECTION 2: REFERENCES
%% ============================================================
\section{References}
\label{sec:references}

\begin{itemize}
  \item \citet[Chapters 1--3]{axler2015} -- Primary reference; develops vector
        spaces, span, independence, bases, and linear maps without determinants.
  \item \citet[Chapters 1--3]{hoffman1971} -- Classical treatment with emphasis on
        the interplay between abstract vector spaces and matrices.
  \item \citet[Chapters 1--3]{strang2016} -- Computational perspective with
        many worked examples connecting to applications.
  \item \citet[Chapter 1]{horn2013} -- Concise review of vector spaces and linear
        maps as background for matrix analysis.
  \item \citet[Chapters 1--4]{lax2007} -- Elegant proofs and connections to
        functional analysis.
\end{itemize}

%% ============================================================
%% SECTION 3: VECTOR SPACES
%% ============================================================
\section{Vector Spaces}
\label{sec:vector_spaces}

Throughout this lesson, $\F$ denotes either $\R$ or $\C$.

\subsection{Axioms}

\begin{definition}[Vector Space]
\label{def:vector_space}
A \emph{vector space} over a field $\F$ is a set $V$ together with an addition
$+\colon V \times V \to V$ and a scalar multiplication $\cdot\colon \F \times V \to V$
satisfying the following axioms for all $u, v, w \in V$ and all $a, b \in \F$:
\begin{enumerate}[nosep]
  \item \textbf{Commutativity:} $u + v = v + u$.
  \item \textbf{Associativity of addition:} $(u + v) + w = u + (v + w)$.
  \item \textbf{Additive identity:} There exists $0 \in V$ such that $v + 0 = v$
        for all $v \in V$.
  \item \textbf{Additive inverse:} For each $v \in V$, there exists $w \in V$
        with $v + w = 0$.
  \item \textbf{Multiplicative identity:} $1v = v$.
  \item \textbf{Associativity of scalar multiplication:} $a(bv) = (ab)v$.
  \item \textbf{Distributive laws:} $a(u + v) = au + av$ and $(a + b)v = av + bv$.
\end{enumerate}
Elements of $V$ are called \emph{vectors} and elements of $\F$ are called \emph{scalars}.
\end{definition}

\begin{remark}
\label{rem:uniqueness_zero}
The additive identity is unique: if $0'$ also satisfies $v + 0' = v$ for all $v$,
then $0 = 0 + 0' = 0'$.  Similarly, each vector has a unique additive inverse,
which we denote $-v$.
\end{remark}

\begin{example}[Standard Vector Spaces]
\label{ex:standard_spaces}
The following are vector spaces over $\F$:
\begin{enumerate}[nosep]
  \item $\F^n = \{(x_1, \ldots, x_n) : x_i \in \F\}$ with componentwise addition
        and scalar multiplication.
  \item $\F[x]_{\leq n}$, the set of polynomials with coefficients in $\F$ of
        degree at most $n$, with usual polynomial operations.
  \item $C[a,b]$, the set of continuous functions $f\colon [a,b] \to \R$, with
        pointwise addition $(f+g)(x) = f(x) + g(x)$ and scalar multiplication
        $(\alpha f)(x) = \alpha f(x)$.
  \item $\F^{m \times n}$, the set of $m \times n$ matrices with entries in $\F$,
        with entry-wise addition and scalar multiplication.
\end{enumerate}
\end{example}

\begin{example}[Solution Set of a Homogeneous System]
\label{ex:solution_space}
Let $A \in \R^{m \times n}$.  The solution set
$S = \{x \in \R^n : Ax = 0\}$ is a vector space over $\R$ (a subspace of $\R^n$):
if $Ax = 0$ and $Ay = 0$, then $A(x + y) = 0$ and $A(\alpha x) = 0$.
In machine learning, the null space of a feature matrix determines the set of
weight vectors that produce identical predictions.
\end{example}

\subsection{Subspaces}

\begin{definition}[Subspace]
\label{def:subspace}
A subset $U \subseteq V$ is a \emph{subspace} of $V$ if $U$ is itself a vector
space under the operations inherited from $V$.  Equivalently (the \emph{subspace
test}), $U$ is a subspace if and only if:
\begin{enumerate}[nosep]
  \item $0 \in U$,
  \item $u + v \in U$ for all $u, v \in U$ \quad (closure under addition),
  \item $\lambda u \in U$ for all $\lambda \in \F$, $u \in U$
        \quad (closure under scalar multiplication).
\end{enumerate}
\end{definition}

\begin{proposition}[Intersection of Subspaces]
\label{prop:intersection}
If $U_1, U_2$ are subspaces of $V$, then $U_1 \cap U_2$ is a subspace of $V$.
More generally, the intersection of any collection of subspaces is a subspace.
\end{proposition}

\begin{proof}
We verify the three conditions of Definition~\ref{def:subspace}.
(1)~Since $0 \in U_1$ and $0 \in U_2$, we have $0 \in U_1 \cap U_2$.
(2)~If $u, v \in U_1 \cap U_2$, then $u, v \in U_1$, so $u + v \in U_1$ by
closure in $U_1$.  Similarly $u + v \in U_2$.  Hence $u + v \in U_1 \cap U_2$.
(3)~If $u \in U_1 \cap U_2$ and $\lambda \in \F$, then $\lambda u \in U_1$ and
$\lambda u \in U_2$, so $\lambda u \in U_1 \cap U_2$.
The argument for an arbitrary collection is identical, checking each subspace
in the family.
\end{proof}

\begin{warning}[title={Union of Subspaces}]
The \emph{union} of two subspaces is generally \emph{not} a subspace.  For instance,
in $\R^2$ let $U_1 = \{(x,0) : x \in \R\}$ and $U_2 = \{(0,y) : y \in \R\}$.
Then $(1,0) \in U_1$ and $(0,1) \in U_2$, but $(1,1) = (1,0) + (0,1) \notin
U_1 \cup U_2$.  The correct notion that combines two subspaces into a larger
subspace is the \emph{sum}, defined next.
\end{warning}

\subsection{Sums and Direct Sums}

\begin{definition}[Sum of Subspaces]
\label{def:sum}
If $U_1, \ldots, U_m$ are subspaces of $V$, their \emph{sum} is
\[
  U_1 + \cdots + U_m
  = \{u_1 + \cdots + u_m : u_i \in U_i \text{ for each } i\}.
\]
\end{definition}

\begin{proposition}
\label{prop:sum_smallest}
The sum $U_1 + \cdots + U_m$ is the smallest subspace of $V$ containing every $U_i$.
\end{proposition}

\begin{proof}
The sum is a subspace: it contains $0$ (take each $u_i = 0$); if
$u = u_1 + \cdots + u_m$ and $v = v_1 + \cdots + v_m$ with $u_i, v_i \in U_i$,
then $u + v = (u_1 + v_1) + \cdots + (u_m + v_m)$ with $u_i + v_i \in U_i$;
and $\lambda u = \lambda u_1 + \cdots + \lambda u_m$ with $\lambda u_i \in U_i$.
Each $U_i$ is contained in the sum (set all other summands to $0$).  Any subspace
containing every $U_i$ must contain all sums $u_1 + \cdots + u_m$, so the sum
is the smallest such subspace.
\end{proof}

\begin{definition}[Direct Sum]
\label{def:direct_sum}
The sum $U_1 + \cdots + U_m$ is called a \emph{direct sum}, written
$U_1 \oplus \cdots \oplus U_m$, if each element of the sum can be expressed
in \emph{exactly one way} as $u_1 + \cdots + u_m$ with $u_i \in U_i$.
\end{definition}

\begin{proposition}[Characterization of Direct Sums for Two Subspaces]
\label{prop:direct_sum_char}
$V = U_1 \oplus U_2$ if and only if $V = U_1 + U_2$ and $U_1 \cap U_2 = \{0\}$.
\end{proposition}

\begin{proof}
$(\Rightarrow)$\; Suppose $V = U_1 \oplus U_2$ and let $v \in U_1 \cap U_2$.
Then $0 = v + (-v)$ with $v \in U_1$ and $-v \in U_2$.  But also $0 = 0 + 0$
with $0 \in U_1$ and $0 \in U_2$.  By the uniqueness requirement in the direct
sum, $v = 0$.

$(\Leftarrow)$\; Suppose $V = U_1 + U_2$ and $U_1 \cap U_2 = \{0\}$.  If
$u_1 + u_2 = u_1' + u_2'$ with $u_i, u_i' \in U_i$, then
$u_1 - u_1' = u_2' - u_2$.  The left side belongs to $U_1$ and the right to
$U_2$, so both belong to $U_1 \cap U_2 = \{0\}$.  Hence $u_1 = u_1'$ and
$u_2 = u_2'$, giving uniqueness.
\end{proof}

\begin{example}[Direct Sum in $\R^3$]
\label{ex:direct_sum_R3}
Let $U = \{(x, y, 0) : x, y \in \R\}$ and $W = \{(0, 0, z) : z \in \R\}$.
Every $(x,y,z) \in \R^3$ has the unique decomposition $(x,y,0) + (0,0,z)$
with the first summand in $U$ and the second in $W$.  Since $U \cap W = \{0\}$,
Proposition~\ref{prop:direct_sum_char} confirms $\R^3 = U \oplus W$.
\end{example}

\begin{theorem}[Dimension of a Sum]
\label{thm:dim_sum}
If $U_1$ and $U_2$ are finite-dimensional subspaces of $V$, then
\[
  \dim(U_1 + U_2) = \dim U_1 + \dim U_2 - \dim(U_1 \cap U_2).
\]
\end{theorem}

\begin{proof}
Let $\{w_1, \ldots, w_k\}$ be a basis of $U_1 \cap U_2$.  Extend this to a
basis $\{w_1, \ldots, w_k, u_1, \ldots, u_p\}$ of $U_1$ (by
Proposition~\ref{prop:basis_extension} below) and to a basis
$\{w_1, \ldots, w_k, v_1, \ldots, v_q\}$ of $U_2$.

We claim the list
$\mathcal{L} = \{w_1, \ldots, w_k, u_1, \ldots, u_p, v_1, \ldots, v_q\}$
is a basis of $U_1 + U_2$.

\textbf{Spanning.}\; Any $x \in U_1 + U_2$ is $x = a + b$ with $a \in U_1$
and $b \in U_2$.  Then $a$ is a combination of the $w$'s and $u$'s, and $b$
is a combination of the $w$'s and $v$'s, so $x$ lies in $\Span(\mathcal{L})$.

\textbf{Independence.}\; Suppose
$\sum_{i} \alpha_i w_i + \sum_{j} \beta_j u_j + \sum_{\ell} \gamma_\ell v_\ell = 0$.
Then $\sum_\ell \gamma_\ell v_\ell = -\sum_i \alpha_i w_i - \sum_j \beta_j u_j$.
The left side lies in $U_2$ and the right side lies in $U_1$, so both sides
belong to $U_1 \cap U_2$.  Hence $\sum_\ell \gamma_\ell v_\ell
= \sum_i \delta_i w_i$ for some scalars $\delta_i$.  Since
$\{w_1, \ldots, w_k, v_1, \ldots, v_q\}$ is a basis of $U_2$ (hence linearly
independent), all $\gamma_\ell = 0$ and all $\delta_i = 0$.  The remaining
equation $\sum_i \alpha_i w_i + \sum_j \beta_j u_j = 0$, together with the
independence of the basis of $U_1$, gives all $\alpha_i = 0$ and $\beta_j = 0$.

Therefore $\dim(U_1 + U_2) = k + p + q = (k+p) + (k+q) - k
= \dim U_1 + \dim U_2 - \dim(U_1 \cap U_2)$.
\end{proof}

\begin{rlconnection}[title={Feature Subspaces in Representation Learning}]
In multi-task learning, each task may learn features in a subspace $U_i$ of
a shared representation space.  The dimension formula
$\dim(U_1 + U_2) = \dim U_1 + \dim U_2 - \dim(U_1 \cap U_2)$ quantifies
how much the tasks share: a large intersection means the tasks exploit
similar features, justifying a shared hidden layer with dimension close
to $\dim U_1$ rather than $\dim U_1 + \dim U_2$.
\end{rlconnection}

With the vocabulary of subspaces and sums in place, we now turn to the
building blocks from which every finite-dimensional subspace is constructed:
spanning sets and linearly independent lists.

%% ============================================================
%% SECTION 4: BASES AND DIMENSION
%% ============================================================
\section{Bases and Dimension}
\label{sec:bases_dimension}

\subsection{Span}

\begin{definition}[Span]
\label{def:span}
The \emph{span} (or \emph{linear span}) of a list of vectors
$v_1, \ldots, v_m \in V$ is
\[
  \Span(v_1, \ldots, v_m)
  = \bigl\{\lambda_1 v_1 + \cdots + \lambda_m v_m
    : \lambda_1, \ldots, \lambda_m \in \F\bigr\}.
\]
By convention, $\Span(\,) = \{0\}$ (the span of the empty list is $\{0\}$).
A vector space $V$ is \emph{finite-dimensional} if there exists a finite list
of vectors that spans $V$.
\end{definition}

\begin{example}[Span in $\R^3$]
\label{ex:span_R3}
The vectors $e_1 = (1,0,0)$ and $e_2 = (0,1,0)$ satisfy
$\Span(e_1, e_2) = \{(a, b, 0) : a, b \in \R\}$, which is the $xy$-plane
in $\R^3$.  Adding $e_3 = (0,0,1)$ gives $\Span(e_1, e_2, e_3) = \R^3$.
\end{example}

\subsection{Linear Independence}

\begin{definition}[Linear Independence]
\label{def:linear_independence}
Vectors $v_1, \ldots, v_m \in V$ are \emph{linearly independent} if the only
choice of scalars satisfying
$\lambda_1 v_1 + \cdots + \lambda_m v_m = 0$
is $\lambda_1 = \cdots = \lambda_m = 0$.  They are \emph{linearly dependent}
otherwise, i.e., some non-trivial combination equals $0$.
\end{definition}

\begin{example}[Testing Independence]
\label{ex:independence_test}
In $\R^3$, the vectors $(1,0,0)$, $(1,1,0)$, $(1,1,1)$ are linearly
independent: if $\alpha(1,0,0) + \beta(1,1,0) + \gamma(1,1,1) = (0,0,0)$,
then the system $\alpha + \beta + \gamma = 0$, $\beta + \gamma = 0$,
$\gamma = 0$ gives $\gamma = 0$, $\beta = 0$, $\alpha = 0$.

On the other hand, $(1,0,0)$, $(0,1,0)$, $(1,1,0)$ are linearly dependent
because $(1,0,0) + (0,1,0) - (1,1,0) = (0,0,0)$.
\end{example}

\begin{lemma}[Linear Dependence Lemma]
\label{lem:linear_dependence}
If $v_1, \ldots, v_m$ are linearly dependent and $v_1 \neq 0$, then there
exists an index $j \in \{2, \ldots, m\}$ such that:
\begin{enumerate}[nosep]
  \item $v_j \in \Span(v_1, \ldots, v_{j-1})$, and
  \item $\Span(v_1, \ldots, v_{j-1}, v_{j+1}, \ldots, v_m)
         = \Span(v_1, \ldots, v_m)$.
\end{enumerate}
In words, one can remove $v_j$ without shrinking the span.
\end{lemma}

\begin{proof}
Since the list is linearly dependent, there exist scalars
$\lambda_1, \ldots, \lambda_m$, not all zero, with
$\sum_{i=1}^{m} \lambda_i v_i = 0$.  Let $j$ be the largest index such that
$\lambda_j \neq 0$.  We must have $j \geq 2$: if $j = 1$, then
$\lambda_1 v_1 = 0$ with $\lambda_1 \neq 0$, giving $v_1 = 0$, contradicting
$v_1 \neq 0$.

Solving for $v_j$:
\[
  v_j = -\lambda_j^{-1} \sum_{i=1}^{j-1} \lambda_i v_i
  \in \Span(v_1, \ldots, v_{j-1}).
\]
For (2), any vector in $\Span(v_1, \ldots, v_m)$ can replace its $v_j$
component using the expression above, so removing $v_j$ does not change
the span.
\end{proof}

\subsection{Basis and Dimension}

\begin{definition}[Basis]
\label{def:basis}
A \emph{basis} of $V$ is a list of vectors in $V$ that is both linearly
independent and spans $V$.
\end{definition}

\begin{example}[Standard Basis of $\F^n$]
\label{ex:standard_basis}
The \emph{standard basis} of $\F^n$ is $e_1 = (1,0,\ldots,0)$,
$e_2 = (0,1,\ldots,0)$, \ldots, $e_n = (0,0,\ldots,1)$.  Independence is
immediate, and any $(x_1, \ldots, x_n) = x_1 e_1 + \cdots + x_n e_n$.
\end{example}

\begin{example}[Basis of Polynomial Space]
\label{ex:polynomial_basis}
A basis of $\F[x]_{\leq n}$ is $\{1, x, x^2, \ldots, x^n\}$.  Thus
$\dim \F[x]_{\leq n} = n + 1$.
\end{example}

\begin{theorem}[All Bases Have the Same Length]
\label{thm:bases_same_length}
If $V$ is finite-dimensional, then any two bases of $V$ have the same number
of elements.
\end{theorem}

\begin{proof}
Let $\mathcal{U} = \{u_1, \ldots, u_m\}$ and $\mathcal{W} = \{w_1, \ldots, w_n\}$
be two bases of $V$.  We show $m \leq n$.

Since $\mathcal{W}$ spans $V$, the list $w_1, \ldots, w_n, u_1$ is linearly
dependent (it has $n+1$ vectors in a space spanned by $n$ vectors).
By the linear dependence lemma (Lemma~\ref{lem:linear_dependence}), since
$w_1 \neq 0$, we can remove some $w_j$ without changing the span.  We thus
obtain a spanning list of length $n$ that includes $u_1$.

Repeat: insert $u_2$ into this list (making it linearly dependent, as it has
$n+1$ elements spanning an $n$-dimensional space) and remove one of the
remaining $w$'s.  After $m$ steps, we have a spanning list consisting of all
$u_1, \ldots, u_m$ together with $n - m$ of the original $w$'s.  This requires
$n - m \geq 0$, so $m \leq n$.

Exchanging the roles of $\mathcal{U}$ and $\mathcal{W}$ gives $n \leq m$.
Therefore $m = n$.
\end{proof}

\begin{keyresult}[title={Dimension Is Well-Defined}]
Theorem~\ref{thm:bases_same_length} ensures that the following definition is
unambiguous: the dimension of a finite-dimensional vector space is the common
length of all its bases.  This single number encodes the ``size'' of the space
in a coordinate-free way.
\end{keyresult}

\begin{definition}[Dimension]
\label{def:dimension}
The \emph{dimension} of a finite-dimensional vector space $V$, denoted
$\dim V$, is the number of elements in any (hence every) basis of $V$.
\end{definition}

\begin{proposition}[Basis Extension]
\label{prop:basis_extension}
Every linearly independent list in a finite-dimensional vector space $V$ can be
extended to a basis of $V$.
\end{proposition}

\begin{proof}
Let $v_1, \ldots, v_m$ be linearly independent in $V$ with $\dim V = n$.
If $m = n$, the list already spans $V$: a linearly independent list of length
equal to the dimension must be a basis (any vector not in the span would extend
it to an independent list of length $n+1$, contradicting
Theorem~\ref{thm:bases_same_length}).

If $m < n$, the list does not span $V$, so there exists
$v_{m+1} \notin \Span(v_1, \ldots, v_m)$.  Then $v_1, \ldots, v_{m+1}$ is
linearly independent (any dependence relation involving $v_{m+1}$ would place
$v_{m+1}$ in $\Span(v_1, \ldots, v_m)$).  Repeat until the list has length $n$.
\end{proof}

\begin{proposition}[Dimension of a Subspace]
\label{prop:dim_subspace}
If $U$ is a subspace of a finite-dimensional space $V$, then
$\dim U \leq \dim V$.  Equality holds if and only if $U = V$.
\end{proposition}

\begin{proof}
Any linearly independent list in $U$ is also linearly independent in $V$,
so its length is at most $\dim V$.  Hence $\dim U \leq \dim V$.

If $\dim U = \dim V = n$, let $\{u_1, \ldots, u_n\}$ be a basis of $U$.
This is a linearly independent list of length $n$ in $V$, so by
Theorem~\ref{thm:bases_same_length} it is a basis of $V$.  Thus every
$v \in V$ lies in $\Span(u_1, \ldots, u_n) \subseteq U$, giving $V \subseteq U$
and hence $U = V$.

Conversely, $U = V$ trivially implies $\dim U = \dim V$.
\end{proof}

\begin{rlconnection}[title={Embedding Dimension and Feature Rank}]
In neural network design, the hidden layer dimension $d$ in an encoder
$\phi\colon \R^n \to \R^d$ controls the representational capacity.
The image $\phi(\R^n)$ (under a linear encoder) is a subspace of $\R^d$ with
$\dim(\phi(\R^n)) = \rank(\phi) \leq \min(n, d)$.  Choosing $d$ too small
forces $\rank(\phi) < n$, losing information.  The dimension inequality
in Proposition~\ref{prop:dim_subspace} makes this precise: the encoder
image, as a subspace, cannot exceed the ambient dimension $d$.
\end{rlconnection}

Having established that finite-dimensional vector spaces are characterized up to
isomorphism by a single integer (the dimension), we now study the maps that
respect vector space structure.

%% ============================================================
%% SECTION 5: LINEAR MAPS
%% ============================================================
\section{Linear Maps}
\label{sec:linear_maps}

\subsection{Definition and Basic Properties}

\begin{definition}[Linear Map]
\label{def:linear_map}
A function $T\colon V \to W$ between vector spaces over $\F$ is a \emph{linear
map} if for all $u, v \in V$ and $\lambda \in \F$:
\begin{enumerate}[nosep]
  \item $T(u + v) = Tu + Tv$ \quad (additivity),
  \item $T(\lambda v) = \lambda\, Tv$ \quad (homogeneity).
\end{enumerate}
Equivalently, $T(\alpha u + \beta v) = \alpha\, Tu + \beta\, Tv$ for all
$\alpha, \beta \in \F$ and $u, v \in V$.  The set of all linear maps from
$V$ to $W$ is denoted $\mathcal{L}(V, W)$.
\end{definition}

\begin{example}[Linear Maps]
\label{ex:linear_maps}
\begin{enumerate}[nosep]
  \item The \emph{zero map} $0\colon V \to W$ defined by $0(v) = 0$ for all $v$.
  \item The \emph{identity map} $I\colon V \to V$ defined by $Iv = v$ for all $v$.
  \item \emph{Differentiation} $D\colon \F[x]_{\leq n} \to \F[x]_{\leq n-1}$
        defined by $Dp = p'$ is linear.
  \item For any matrix $A \in \F^{m \times n}$, the map $T_A\colon \F^n \to \F^m$
        defined by $T_A(x) = Ax$ is linear.
\end{enumerate}
\end{example}

\subsection{Null Space and Range}

\begin{definition}[Null Space and Range]
\label{def:null_range}
For a linear map $T\colon V \to W$:
\begin{itemize}[nosep]
  \item The \emph{null space} (or \emph{kernel}) of $T$ is
        $\nullspace T = \{v \in V : Tv = 0\}$.
  \item The \emph{range} (or \emph{image}) of $T$ is
        $\range T = \{Tv : v \in V\}$.
\end{itemize}
\end{definition}

\begin{proposition}
\label{prop:null_range_subspace}
For any linear map $T\colon V \to W$, the null space $\nullspace T$ is a
subspace of $V$ and the range $\range T$ is a subspace of $W$.
\end{proposition}

\begin{proof}
\textbf{Null space.}\;
$T(0) = 0$, so $0 \in \nullspace T$.  If $u, v \in \nullspace T$, then
$T(u+v) = Tu + Tv = 0 + 0 = 0$, so $u + v \in \nullspace T$.  If
$\lambda \in \F$, then $T(\lambda u) = \lambda\, Tu = \lambda \cdot 0 = 0$,
so $\lambda u \in \nullspace T$.

\textbf{Range.}\;
$T(0) = 0$, so $0 \in \range T$.  If $w_1, w_2 \in \range T$, write
$w_1 = Tv_1$ and $w_2 = Tv_2$.  Then $w_1 + w_2 = T(v_1 + v_2) \in \range T$.
If $\lambda \in \F$, then $\lambda w_1 = T(\lambda v_1) \in \range T$.
\end{proof}

\begin{proposition}[Injectivity Criterion]
\label{prop:injectivity}
A linear map $T\colon V \to W$ is injective if and only if
$\nullspace T = \{0\}$.
\end{proposition}

\begin{proof}
$(\Rightarrow)$\; If $T$ is injective and $Tv = 0 = T(0)$, then $v = 0$.
Hence $\nullspace T = \{0\}$.

$(\Leftarrow)$\; Suppose $\nullspace T = \{0\}$ and $Tu = Tv$.  Then
$T(u - v) = Tu - Tv = 0$, so $u - v \in \nullspace T = \{0\}$, giving
$u = v$.
\end{proof}

\subsection{The Rank--Nullity Theorem}

\begin{keyresult}[title={The Rank--Nullity Theorem}]
\begin{theorem}[Rank--Nullity Theorem]
\label{thm:rank_nullity}
If $V$ is finite-dimensional and $T\colon V \to W$ is a linear map, then
\[
  \dim V = \dim(\nullspace T) + \dim(\range T).
\]
Equivalently, writing $\nullity(T) = \dim(\nullspace T)$ and
$\rank(T) = \dim(\range T)$:
\[
  \dim V = \nullity(T) + \rank(T).
\]
\end{theorem}
\end{keyresult}

\begin{proof}
Let $\{u_1, \ldots, u_m\}$ be a basis of $\nullspace T$.  By
Proposition~\ref{prop:basis_extension}, extend it to a basis
$\{u_1, \ldots, u_m, v_1, \ldots, v_k\}$ of $V$, so that
$\dim V = m + k$.  We claim that $\{Tv_1, \ldots, Tv_k\}$ is a basis
of $\range T$.

\textbf{Spanning.}\;
Let $w \in \range T$.  Then $w = T(x)$ for some $x \in V$.  Write
$x = \alpha_1 u_1 + \cdots + \alpha_m u_m + \mu_1 v_1 + \cdots + \mu_k v_k$.
Applying $T$ and using $Tu_i = 0$:
\[
  w = T(x) = \mu_1 Tv_1 + \cdots + \mu_k Tv_k.
\]
Hence $\{Tv_1, \ldots, Tv_k\}$ spans $\range T$.

\textbf{Independence.}\;
Suppose $\mu_1 Tv_1 + \cdots + \mu_k Tv_k = 0$.  By linearity,
$T(\mu_1 v_1 + \cdots + \mu_k v_k) = 0$, so
$\mu_1 v_1 + \cdots + \mu_k v_k \in \nullspace T$.  Since
$\{u_1, \ldots, u_m\}$ is a basis of $\nullspace T$, there exist scalars
$\alpha_i$ with
\[
  \mu_1 v_1 + \cdots + \mu_k v_k = \alpha_1 u_1 + \cdots + \alpha_m u_m.
\]
Rearranging:
$\alpha_1 u_1 + \cdots + \alpha_m u_m - \mu_1 v_1 - \cdots - \mu_k v_k = 0$.
Since $\{u_1, \ldots, u_m, v_1, \ldots, v_k\}$ is a basis of $V$ (hence
linearly independent), all coefficients vanish: $\alpha_i = 0$ for each $i$
and $\mu_j = 0$ for each $j$.

Therefore $\dim(\range T) = k$ and
$\dim V = m + k = \dim(\nullspace T) + \dim(\range T)$.
\end{proof}

\begin{corollary}[Equivalence for Square Dimensions]
\label{cor:square}
Let $T\colon V \to W$ be linear with $\dim V = \dim W = n$.  The following
are equivalent:
\begin{enumerate}[nosep,label=(\roman*)]
  \item $T$ is injective.
  \item $T$ is surjective.
  \item $T$ is invertible (bijective).
\end{enumerate}
\end{corollary}

\begin{proof}
By the rank--nullity theorem, $\nullity(T) + \rank(T) = n$.

(i)$\Rightarrow$(ii):\; Injectivity means $\nullity(T) = 0$, so
$\rank(T) = n = \dim W$, i.e., $\range T = W$ (surjectivity).

(ii)$\Rightarrow$(i):\; Surjectivity means $\rank(T) = n$, so
$\nullity(T) = 0$ (injectivity).

(iii)$\Leftrightarrow$(i)$\wedge$(ii):\; Invertibility is, by definition,
injectivity plus surjectivity, and we have shown these are equivalent.
\end{proof}

\begin{example}[Rank--Nullity for Differentiation]
\label{ex:rank_nullity_diff}
Let $D\colon \R[x]_{\leq 3} \to \R[x]_{\leq 2}$ be the differentiation map
$Dp = p'$.  Here $\dim(\R[x]_{\leq 3}) = 4$ and
$\nullspace D = \Span(1)$ (the constant polynomials), so
$\nullity(D) = 1$.  By rank--nullity, $\rank(D) = 4 - 1 = 3 = \dim(\R[x]_{\leq 2})$,
confirming that $D$ is surjective (every polynomial of degree $\leq 2$ is the
derivative of some cubic).
\end{example}

\begin{example}[Projection Map]
\label{ex:projection}
Define $P\colon \R^3 \to \R^3$ by $P(x_1, x_2, x_3) = (x_1, x_2, 0)$.
Then $\nullspace P = \{(0,0,z) : z \in \R\}$ with $\nullity(P) = 1$,
and $\range P = \{(x,y,0) : x,y \in \R\}$ with $\rank(P) = 2$.
Indeed $1 + 2 = 3 = \dim \R^3$.
\end{example}

\begin{rlconnection}[title={Rank Deficiency and Bottleneck Layers}]
In a deep neural network, each linear layer $W_\ell \in \R^{d_{\ell} \times d_{\ell-1}}$
defines a linear map $\R^{d_{\ell-1}} \to \R^{d_\ell}$.  By rank--nullity,
$\rank(W_\ell) + \nullity(W_\ell) = d_{\ell-1}$.  If $d_\ell < d_{\ell-1}$
(a bottleneck), then $\nullity(W_\ell) \geq d_{\ell-1} - d_\ell > 0$:
information is inevitably lost.  This is precisely the mechanism exploited
in autoencoders, where the bottleneck forces the network to learn a
low-dimensional representation.
\end{rlconnection}

\subsection{Matrix Representation}

Fix ordered bases $\mathcal{B} = \{v_1, \ldots, v_n\}$ of $V$ and
$\mathcal{C} = \{w_1, \ldots, w_m\}$ of $W$.  A linear map $T\colon V \to W$
is uniquely determined by its action on the basis of $V$.  Write
\[
  Tv_j = \sum_{i=1}^{m} A_{ij}\, w_i, \qquad j = 1, \ldots, n.
\]
The $m \times n$ matrix $A = (A_{ij})$ is called the \emph{matrix of $T$ with
respect to $\mathcal{B}$ and $\mathcal{C}$}, denoted
$\mathcal{M}(T, \mathcal{B}, \mathcal{C})$ or simply $[T]_{\mathcal{B}}^{\mathcal{C}}$.

\begin{proposition}
\label{prop:composition_matrix}
If $S\colon U \to V$ has matrix $B$ (with respect to appropriate bases) and
$T\colon V \to W$ has matrix $A$, then $T \circ S$ has matrix $AB$.
\end{proposition}

\begin{proof}
Let $\{e_1, \ldots, e_p\}$, $\{v_1, \ldots, v_n\}$, $\{w_1, \ldots, w_m\}$ be
the ordered bases of $U$, $V$, $W$ respectively.  Then
$Se_k = \sum_j B_{jk} v_j$ and $Tv_j = \sum_i A_{ij} w_i$.  Hence
\[
  (T \circ S)(e_k)
  = T\Bigl(\sum_j B_{jk} v_j\Bigr)
  = \sum_j B_{jk} Tv_j
  = \sum_j B_{jk} \sum_i A_{ij} w_i
  = \sum_i \Bigl(\sum_j A_{ij} B_{jk}\Bigr) w_i
  = \sum_i (AB)_{ik}\, w_i.
\]
So the matrix of $T \circ S$ is $AB$.
\end{proof}

\begin{example}[Matrix of Differentiation]
\label{ex:matrix_diff}
Let $D\colon \R[x]_{\leq 2} \to \R[x]_{\leq 1}$ be differentiation.  With
basis $\mathcal{B} = \{1, x, x^2\}$ for the domain and
$\mathcal{C} = \{1, x\}$ for the codomain:
$D(1) = 0 = 0 \cdot 1 + 0 \cdot x$,\;
$D(x) = 1 = 1 \cdot 1 + 0 \cdot x$,\;
$D(x^2) = 2x = 0 \cdot 1 + 2 \cdot x$.
Thus
\[
  [D]_{\mathcal{B}}^{\mathcal{C}}
  = \begin{pmatrix} 0 & 1 & 0 \\ 0 & 0 & 2 \end{pmatrix}.
\]
\end{example}

\subsection{Change of Basis and Similarity}

\begin{definition}[Change-of-Basis Matrix]
\label{def:change_of_basis}
Let $\mathcal{B} = \{v_1, \ldots, v_n\}$ and
$\mathcal{B}' = \{v_1', \ldots, v_n'\}$ be two ordered bases of $V$.
The \emph{change-of-basis matrix} from $\mathcal{B}$ to $\mathcal{B}'$ is the
$n \times n$ matrix $P$ whose $j$-th column is the coordinate vector of $v_j$
with respect to $\mathcal{B}'$:
\[
  v_j = \sum_{i=1}^{n} P_{ij}\, v_i', \qquad j = 1, \ldots, n.
\]
Equivalently, if $[x]_{\mathcal{B}}$ denotes the coordinate vector of $x$
relative to $\mathcal{B}$, then $[x]_{\mathcal{B}} = P\,[x]_{\mathcal{B}'}$
for all $x \in V$... wait, let us be precise.  We have
$P\,[x]_{\mathcal{B}'} = [x]_{\mathcal{B}}$, i.e., $P$ converts
$\mathcal{B}'$-coordinates to $\mathcal{B}$-coordinates.
\end{definition}

\begin{remark}
The matrix $P$ is invertible since it represents the identity map
$I\colon V \to V$ with respect to the bases $\mathcal{B}'$ (domain) and
$\mathcal{B}$ (codomain).  Its inverse $P^{-1}$ converts
$\mathcal{B}$-coordinates to $\mathcal{B}'$-coordinates.
\end{remark}

\begin{proposition}[Similarity of Matrix Representations]
\label{prop:similarity}
Let $T\colon V \to V$ be a linear operator.  If $T$ has matrix $A$ with respect
to basis $\mathcal{B}$ and matrix $A'$ with respect to basis $\mathcal{B}'$,
then
\[
  A' = P^{-1} A P,
\]
where $P$ is the change-of-basis matrix from $\mathcal{B}'$ to $\mathcal{B}$
(so that $P[x]_{\mathcal{B}'} = [x]_{\mathcal{B}}$).
\end{proposition}

\begin{proof}
For any $x \in V$, coordinate representations satisfy:
\begin{align*}
  A'[x]_{\mathcal{B}'}
  &= [Tx]_{\mathcal{B}'}
   = P^{-1}[Tx]_{\mathcal{B}}
   = P^{-1} A\, [x]_{\mathcal{B}}
   = P^{-1} A\, P\, [x]_{\mathcal{B}'}.
\end{align*}
Since this holds for all $x \in V$ and the coordinate map relative to
$\mathcal{B}'$ is an isomorphism, $A' = P^{-1}AP$.
\end{proof}

\begin{example}[Change of Basis in $\R^2$]
\label{ex:change_basis_R2}
Let $T\colon \R^2 \to \R^2$ be the linear map with matrix
$A = \bigl(\begin{smallmatrix} 2 & 1 \\ 0 & 3 \end{smallmatrix}\bigr)$
relative to the standard basis $\mathcal{B} = \{e_1, e_2\}$.
Consider the new basis $\mathcal{B}' = \{(1,0), (1,1)\}$.  The change-of-basis
matrix from $\mathcal{B}'$ to $\mathcal{B}$ has columns
$[v_1']_{\mathcal{B}} = (1,0)^\top$ and $[v_2']_{\mathcal{B}} = (1,1)^\top$:
\[
  P = \begin{pmatrix} 1 & 1 \\ 0 & 1 \end{pmatrix}, \qquad
  P^{-1} = \begin{pmatrix} 1 & -1 \\ 0 & 1 \end{pmatrix}.
\]
Then
\[
  A' = P^{-1}AP
  = \begin{pmatrix} 1 & -1 \\ 0 & 1 \end{pmatrix}
    \begin{pmatrix} 2 & 1 \\ 0 & 3 \end{pmatrix}
    \begin{pmatrix} 1 & 1 \\ 0 & 1 \end{pmatrix}
  = \begin{pmatrix} 2 & 0 \\ 0 & 3 \end{pmatrix}.
\]
The new basis diagonalizes $T$.  This illustrates the general strategy:
find a basis of eigenvectors to simplify the matrix representation.
\end{example}

\subsection{Isomorphism}

\begin{proposition}[Isomorphism Criterion]
\label{prop:isomorphism}
Two finite-dimensional vector spaces over $\F$ are isomorphic if and only if
they have the same dimension.
\end{proposition}

\begin{proof}
$(\Rightarrow)$\; Let $T\colon V \to W$ be an isomorphism and let
$\{v_1, \ldots, v_n\}$ be a basis of $V$.  We show
$\{Tv_1, \ldots, Tv_n\}$ is a basis of $W$.

\emph{Independence:}\; If $\sum_i \lambda_i Tv_i = 0$, then
$T(\sum_i \lambda_i v_i) = 0$.  Since $T$ is injective,
$\sum_i \lambda_i v_i = 0$, and since $\{v_i\}$ is independent,
all $\lambda_i = 0$.

\emph{Spanning:}\; For any $w \in W$, surjectivity of $T$ gives $w = Tv$
for some $v = \sum_i \lambda_i v_i$.  Then
$w = \sum_i \lambda_i Tv_i \in \Span(Tv_1, \ldots, Tv_n)$.

Hence $\dim W = n = \dim V$.

$(\Leftarrow)$\; Suppose $\dim V = \dim W = n$.  Pick bases
$\{v_1, \ldots, v_n\}$ of $V$ and $\{w_1, \ldots, w_n\}$ of $W$.
Define $T\colon V \to W$ by $Tv_i = w_i$ for each $i$ and extend linearly.
This map is well-defined (every $v \in V$ has a unique basis expansion) and
linear by construction.  It is bijective because it maps a basis to a basis:
the same argument as above shows $\{Tv_1, \ldots, Tv_n\} = \{w_1, \ldots, w_n\}$
is a basis of $W$.
\end{proof}

\begin{rlconnection}[title={Bellman Equation as a Linear System}]
For a fixed policy $\pi$ in a finite MDP with $n$ states, the Bellman
equation is $V^\pi = r^\pi + \gamma P^\pi V^\pi$, where
$V^\pi, r^\pi \in \R^n$ and $P^\pi \in \R^{n \times n}$ is the transition
matrix.  Rearranging gives $(I - \gamma P^\pi) V^\pi = r^\pi$.  The
operator $I - \gamma P^\pi$ is a linear map $\R^n \to \R^n$.  Since
$\gamma < 1$ and all eigenvalues of $P^\pi$ have modulus at most $1$, the
eigenvalues of $I - \gamma P^\pi$ are $1 - \gamma \lambda_i \neq 0$, so
the map is invertible by Corollary~\ref{cor:square}.  Hence the Bellman
equation has a unique solution.
\end{rlconnection}

%% ============================================================
%% SECTION 6: EXERCISES
%% ============================================================
\section{Exercises}
\label{sec:exercises}

\begin{exercise}[Subspace Verification]
\label{ex:subspace_verify}
\textnormal{(\(*\))} \\
Let $V = \R^3$.  Determine which of the following subsets are subspaces.
Prove your answer in each case.
\begin{enumerate}[nosep,label=(\alph*)]
  \item $S_1 = \{(x,y,z) \in \R^3 : x + 2y - z = 0\}$.
  \item $S_2 = \{(x,y,z) \in \R^3 : x + 2y - z = 1\}$.
  \item $S_3 = \{(x,y,z) \in \R^3 : xz = 0\}$.
\end{enumerate}
\end{exercise}

\begin{exercise}[Dimension of a Sum]
\label{ex:dim_sum}
\textnormal{(\(*\))} \\
Let $U = \Span\bigl((1,1,0,0),\, (1,0,1,0)\bigr)$ and
$W = \Span\bigl((1,0,1,0),\, (0,0,1,1)\bigr)$ in $\R^4$.
Compute $\dim(U \cap W)$, $\dim(U + W)$, and determine whether
$U + W$ is a direct sum.
\end{exercise}

\begin{exercise}[Rank--Nullity Application]
\label{ex:rank_nullity_app}
\textnormal{(\(**\))} \\
Let $T\colon \R^4 \to \R^3$ be the linear map defined by
\[
  T(x_1, x_2, x_3, x_4) = (x_1 + x_2,\; x_2 + x_3,\; x_3 + x_4).
\]
\begin{enumerate}[nosep,label=(\alph*)]
  \item Find a basis for $\nullspace T$ and compute $\nullity(T)$.
  \item Find a basis for $\range T$ and compute $\rank(T)$.
  \item Verify the rank--nullity theorem for this map.
\end{enumerate}
\end{exercise}

\begin{exercise}[Change of Basis and Diagonalization]
\label{ex:change_basis}
\textnormal{(\(**\))} \\
Let $T\colon \R^2 \to \R^2$ have matrix
$A = \bigl(\begin{smallmatrix} 5 & 3 \\ -6 & -4 \end{smallmatrix}\bigr)$
with respect to the standard basis.
\begin{enumerate}[nosep,label=(\alph*)]
  \item Find the eigenvalues of $A$.
  \item For each eigenvalue, find an eigenvector and verify it forms a basis
        of $\R^2$.
  \item Compute the change-of-basis matrix $P$ and verify
        $P^{-1}AP$ is diagonal.
\end{enumerate}
\textit{Hint:} The characteristic polynomial is $\lambda^2 - \lambda - 2$.
\end{exercise}

\begin{exercise}[RL Application: Feature Matrix Rank]
\label{ex:feature_rank}
\textnormal{(\(**\))} \\
In linear value function approximation, the value of state $s_i$ is
approximated as $\hat{V}(s_i) = \phi(s_i)^\top w$ where
$\phi(s_i) \in \R^d$ is a feature vector.  For $n$ states, define the
feature matrix $\Phi \in \R^{n \times d}$ with row $i$ equal to
$\phi(s_i)^\top$.
\begin{enumerate}[nosep,label=(\alph*)]
  \item Show that $\hat{V} = \Phi w$ defines a linear map
        $T_\Phi\colon \R^d \to \R^n$ and that $\range T_\Phi$ is the set of
        representable value functions.
  \item Prove that if $d < n$ and $\rank(\Phi) = d$, then $T_\Phi$ is
        injective but not surjective: there exist value functions that
        cannot be represented.
  \item What happens if $\rank(\Phi) < d$?  Describe the null space of
        $T_\Phi$ and explain its implications for learning.
\end{enumerate}
\end{exercise}

\begin{exercise}[Direct Sum Decomposition]
\label{ex:direct_sum_proof}
\textnormal{(\(***\))} \\
Let $V$ be a finite-dimensional vector space and $U$ a subspace of $V$.
Prove that there exists a subspace $W$ of $V$ such that $V = U \oplus W$.

\textit{Hint:} Extend a basis of $U$ to a basis of $V$ and let $W$ be the
span of the added vectors.
\end{exercise}

%% ============================================================
%% SECTION 7: SUMMARY
%% ============================================================
\section{Summary}
\label{sec:summary}

\begin{itemize}
  \item A \emph{vector space} over $\F$ is a set with addition and scalar
        multiplication satisfying seven axioms (commutativity, associativity,
        identities, inverses, distributivity).

  \item A \emph{subspace} is a nonempty subset closed under addition and
        scalar multiplication; equivalently, it contains $0$ and is closed
        under linear combinations.  The intersection of subspaces is always
        a subspace, but the union generally is not.

  \item The \emph{sum} $U_1 + U_2$ is the smallest subspace containing both
        $U_1$ and $U_2$; it is a \emph{direct sum} $U_1 \oplus U_2$
        precisely when $U_1 \cap U_2 = \{0\}$.

  \item A \emph{basis} is a linearly independent spanning list.  All bases
        of a finite-dimensional space have the same length, which defines
        the \emph{dimension}.

  \item The \emph{linear dependence lemma} is the engine behind most
        dimension arguments: in a dependent list with nonzero first vector,
        some later vector can be removed without changing the span.

  \item The \emph{rank--nullity theorem} states
        $\dim V = \nullity(T) + \rank(T)$ for any linear map
        $T\colon V \to W$ with $V$ finite-dimensional.  This is the
        fundamental constraint relating injectivity, surjectivity, and
        dimension.

  \item Every linear map between finite-dimensional spaces has a unique
        \emph{matrix representation} once bases are chosen.  Changing the
        basis yields a \emph{similar} matrix $A' = P^{-1}AP$.

  \item Two finite-dimensional vector spaces are \emph{isomorphic} if and
        only if they have the same dimension.
\end{itemize}

%% ============================================================
%% SECTION 8: FURTHER READING
%% ============================================================
\section{Further Reading}
\label{sec:further_reading}

\begin{itemize}
  \item \textbf{Lesson 6} (Eigenvalues and Eigenvectors) builds directly on
        the matrix representation and similarity theory developed here.
        The goal is to find bases in which the matrix of a linear operator
        takes a particularly simple form (diagonal or upper triangular).

  \item \textbf{Lesson 7} (Inner Product Spaces) introduces additional
        geometric structure -- lengths, angles, and orthogonality -- on
        vector spaces.  The key results (Gram--Schmidt, orthogonal projections)
        require the basis and dimension theory from this lesson.

  \item \textbf{Lesson 8} (Spectral Theorem and SVD) combines eigenvalue
        theory with inner product structure to obtain the spectral theorem
        for self-adjoint operators and the singular value decomposition.

  \item For a deeper treatment of infinite-dimensional vector spaces and
        their applications to ML, see Phase~01 on functional analysis
        (metric spaces, normed spaces, Hilbert spaces), which generalizes
        the finite-dimensional theory developed here.

  \item \citet{axler2015} is the primary reference for the ``determinant-free''
        approach to linear algebra followed in this lesson.
        \citet{hoffman1971} provides an alternative, more algebraic perspective.
        \citet{lax2007} connects the finite-dimensional theory to operator
        theory and applications.
\end{itemize}

\bibliography{linear-algebra-refs}
\end{document}
